\chapter{记录、回放与外部工具}
\begin{notebox}
\textbf{目标}\quad 保存一段传感器观测，关闭模拟器后重新生成里程计与地图；学会检查时间和输入一致性。\\
\textbf{先修}\quad 第 06、09、10 章；终端已加载主机 ROS 与课程工作空间。
\end{notebox}
\section{把一次运行变成可以反复使用的输入}
调试 SLAM 时，重新运动所产生的传感器序列可能因初值、运动或噪声而变化，难以判断结果变化来自参数还是输入。
rosbag2 把话题消息和时间存进文件；之后可将同一批扫描和 IMU 测量重新送给估计器。
这样调整匹配阈值或滤波噪声后，可以在相同观测上比较结果。

本实验只录制 \code{/scan} 与 \code{/imu/data_raw}。
估计里程计、map/odom 变换和地图由新运行的算法重新计算；传感器安装外参由回放 launch 提供。
如果同时播放旧估计的 TF，新旧两个算法会发布同一条变换边，显示和坐标查询便可能交替使用不同结果。
因此先画清“保存的输入”和“重新计算的输出”：
\begin{center}\begin{minipage}{\linewidth}\makeatletter\def\@captype{figure}\makeatother
\centering
\begin{tikzpicture}[>=stealth,font=\small,
 box/.style={draw=tealblue,fill=pale,rounded corners,minimum height=.85cm,align=center}]
\node[box,minimum width=2cm] (bag) at (0,0) {传感器 bag};
\node[box,minimum width=2.5cm] (player) at (4,0) {播放器};
\node[box,minimum width=3cm] (est) at (9,0) {雷达里程计与融合\\关键帧 SLAM};
\node[box,minimum width=3cm] (out) at (9,-2) {新里程计、TF\\观测地图与点云};
\draw[->] (bag)--(player);
\draw[->] (player)--node[above]{激光、IMU}(est);
\draw[->] (player.south)--++(0,-.85)-|node[pos=.25,below]{/clock}(est.south west);
\draw[->] (est)--(out);
\end{tikzpicture}
\caption{回放的数据流。文件提供固定观测，播放器提供仿真时钟，估计器从观测重新计算状态与地图。}
\end{minipage}\end{center}

\section{先区分采集、存储与墙钟时间}
一条消息涉及三种时刻：
\begin{center}\small
\begin{tabular}{lp{10cm}}\toprule
符号 & 含义与用途\\\midrule
$t_s$ & 消息 header.stamp 中的采集时间，决定测量对应哪个物理状态\\
$t_r$ & 录制器写入 bag 的接收/存储时间，播放器据此安排发送顺序\\
$w$ & 运行程序所经历的墙钟时间，用于观察播放速度与计算耗时\\\bottomrule
\end{tabular}
\end{center}
使用仿真时间录制时，$t_r$ 取录制器最近收到的 /clock 值；$t_s$ 则已在传感器消息中填写。
两个话题到达的先后可能不同，因此二者可相差一个时钟更新间隔。
例如消息的采集时间为 1.000 s，而录制器刚看到的时钟仍为 0.995 s，文件会保留这两个值。
估计器积分使用相邻 header 的差，存储时间用于重放调度。
具体选项可查阅 \href{https://github.com/ros2/rosbag2#simulation-time}{rosbag2 的仿真时间说明}，
并用本机 \code{ros2 bag record --help} 核对当前命令行。

\subsection{先录制一小段传感器}
按第 10 章启动仿真，保持暂停。在另一个终端启动录制器，再恢复仿真：
\begin{lstlisting}[language=bash,numbers=none]
# Terminal B; choose a fresh output directory.
ros2 bag record --use-sim-time -o tmp/my_sensors \
  --topics /scan /imu/data_raw
# Terminal C, after the recorder discovers both topics:
ros2 service call /sim/pause std_srvs/srv/SetBool "{data: false}"
\end{lstlisting}
录制数秒后在终端 B 按 Ctrl+C，让文件关闭；随后停止仿真与旧估计节点。
录制器收到第一个 /clock 后才开始写入消息，因此最初进入文件的采样可能晚于仿真零时刻。
用下面的命令检查文件实际包含的话题、条数和时间范围：
\begin{lstlisting}[language=bash,numbers=none]
ros2 bag info tmp/my_sensors
\end{lstlisting}
如果条数明显偏少，先检查录制器何时发现传感器、QoS 是否匹配，再检查发布频率。
消息内部的 frame、量程、协方差、噪声样本和采集时间均随消息保存。

\section{半速播放为什么仍用原来的积分步长}
在新的终端中启动只含估计节点的回放栈，再启动播放器：
\begin{lstlisting}[language=bash,numbers=none]
# Terminal A; this launch contains estimators and sensor extrinsics.
ros2 launch motion2d_bringup replay_slam.launch.py rviz:=true
# Terminal B; the simulator and old estimators have stopped.
ros2 bag play tmp/my_sensors --clock 200 --rate 0.5
\end{lstlisting}
播放器是这一回放系统的 /clock 发布者，估计节点均使用仿真时间。
设播放倍率为 $r>0$，文件中两条消息的存储时差为 $\Delta t_r$，
则计划的墙钟发送间隔为 $\Delta w=\Delta t_r/r$。
半速 $r=0.5$ 时，计算机有两倍墙钟时间处理同一段运动；消息内部 $t_s$ 保持原值。

例如陀螺测量为 0.2 rad/s，两条 IMU 采集时间相差 0.005 s，积分得到
\[
 \Delta\psi=\omega\Delta t_s=0.2\times0.005=0.001\ \mathrm{rad}.
\]
半速回放时，这两条消息可能相隔约 0.010 s 墙钟时间到达，角度增量仍为 0.001 rad。
若改用到达间隔积分，就会错误地将机器人转速或运动时长放大。

\begin{center}\begin{minipage}{\linewidth}\makeatletter\def\@captype{figure}\makeatother
\centering
\begin{tikzpicture}[x=1cm,y=1cm,>=stealth,font=\small]
\node[anchor=east] at (-.6,0){采集时刻 / s};
\node[anchor=east] at (-.6,-1.7){回放经过时间 / s};
\draw[->,muted] (0,0)--(8,0);\draw[->,muted] (0,-1.7)--(8,-1.7);
\foreach \x/\s/\w in {1/1.000/0,4/1.005/0.010,7/1.010/0.020}{
 \draw[dashed,linecolor] (\x,-.15)--(\x,-1.55);
 \fill[tealblue] (\x,0) circle (.05) node[above]{$\s$};
 \fill[navy] (\x,-1.7) circle (.05) node[below]{$\w$};
}
\node[fill=white,inner sep=3pt] at (4,-.85){同一条消息，header 保持原值};
\end{tikzpicture}
\caption{半速播放的时间对照。为突出比例，图中假设存储时差等于采集时差；实际回放发送顺序由文件时间决定。}
\end{minipage}\end{center}

\subsection{TF 每条边由谁计算}
回放 launch 复用雷达里程计、IMU 融合、估计里程计适配与关键帧 SLAM 四个节点。
雷达里程计提供扫描位姿，融合节点结合 IMU，适配器将融合结果转换为统一里程计接口。
TF 分工为：
\begin{center}\small
\begin{tabular}{ll}\toprule
变换边 & 发布者\\\midrule
map $\to$ odom & 关键帧 SLAM\\
odom $\to$ base\_link & 估计里程计适配器\\
base\_link $\to$ laser & 静态雷达外参发布器\\
base\_link $\to$ imu\_link & 静态 IMU 外参发布器\\\bottomrule
\end{tabular}
\end{center}
课程传感器与圆盘中心重合、朝向一致，所以两条外参是零平移、零旋转。
更换真实记录时，应填入标定的安装位姿，使扫描与惯性测量转换到正确机体系。
RViz 的 Fixed Frame 设为 map，观察新生成的占据地图、注册点云与估计轨迹。
若点云跟随错误的坐标移动，先检查 frame 名称和这四条边的时间，再调整 SLAM 参数。

\section{用自动实验追踪每一条输入}
在仓库根目录可运行完整录制与回放例子，输出路径应是新的目录：
\begin{lstlisting}[language=bash,numbers=none]
/usr/bin/python3 scripts/run_appD.py --output tmp/my_replay
\end{lstlisting}
脚本在独立 ROS 域内启动暂停的 reference 模型，等待录制器订阅两类传感器，推进至六秒后停止。
它关闭录制器和模拟器，再启动上述估计栈并半速回放。
场景采用 seed42 随机世界、720 束/10 Hz 激光和 200 Hz 白噪声 IMU，参数来自第 10 章。

\subsection{先比较字段，再解释结果}
本次记录中的每条传感器消息都有本话题内唯一的采集时间戳。脚本建立
“话题 $+$ 采集纳秒 $\to$ 字段摘要”的对应，用它配对现场接收、文件读取与回放接收的消息。
纳秒值为
\[
 t_{\rm ns}=10^9\,\text{sec}+\text{nanosec},
\]
保留整数能避免把纳秒转换成浮点秒后丢失末位。
字段摘要的核心是下面几行：
\lstinputlisting[language=Python,linerange={replay_digest_begin-replay_digest_end},includerangemarker=false]{../scripts/run_appD.py}
\code{message_to_ordereddict} 展开消息的每个字段，JSON 固定键顺序与分隔符，
SHA-256 把该表示变为便于比较的固定长度摘要。相同字段将产生相同摘要；发现不同摘要时，
再展开消息查找具体差异。量程中的 NaN、无回波、frame、stamp 和协方差都参与比较。

这样比较关注测量本身。CDR 序列化还可能插入用于对齐的填充字节，重新编码时填充值可以变化。
例如某个按 4 字节对齐的字段前已有 1 字节内容，就需补 3 字节以移动到下一个 4 字节边界。
这些填充用于存放布局，解码后的测量数值由真正字段决定。
因此检查传感器重放时，先解码再比较字段，可以排除填充差异的干扰。

\subsection{怎样读这次记录的统计量}
随教材保存的一组结果为：59 帧扫描、1180 帧 IMU，存储时间从 0.105 到 6.000 s，
存储与采集时刻最大相差 0.005 s。
初始时钟与订阅建立影响了文件起点，因此应以文件实际条数为准。
回放结束后，最后里程计时刻为 6.000 s，地图有 12556 个已观测格，累计点云有 537 点。
它们分别说明输入记录长度、估计处理进度以及地图是否已累积观测。
对应字段见 \code{textbook/data/appD_replay.json}。

先确认回放输入逐字段一致，再检查估计输出的完整性。
研究定位精度时，另保存独立的评价真值，按第 19 章做时间对齐并计算 ATE/RPE；
研究地图质量时，在已观测区域比较墙体位置、厚度与占据分类。
这种顺序有助于区分输入丢失、时间处理错误和算法估计误差。

\section{在同一记录上接入外部二维 SLAM}
选做对照可使用 \href{https://github.com/SteveMacenski/slam_toolbox#api}{slam\_toolbox 的官方接口说明}：
输入扫描与 odom 到 base 的有效变换，由算法输出地图。
先为目标版本整理输入表，写清扫描话题、外参、里程计来源和坐标名称。
如果一个系统直接融合 IMU，另一个通过融合里程计间接使用 IMU，应把这条信息路径写进比较条件。

对照时按如下顺序准备：
\begin{enumerate}
\item 固定同一份传感器 bag 与外参；确定每个算法可获得哪些观测和先验。
\item 串行运行两个系统，使每次只有一个发布者计算 map/odom；每次从新的算法状态开始。
\item 固定回放速率、初值和计算预算，保存参数、版本、状态及输出轨迹。
\item 用同一评价真值、时间交集与 ATE/RPE 定义比较；同时记录失败次数与地图质量。
\end{enumerate}
若某个系统获得更精确的里程计先验，结果就同时反映先验与算法的作用。
可先比较整套系统，再固定共同先验做第二组实验，逐步缩小影响因素。

\section{从教材找到函数说明，再回到代码}
教材按问题展开推导，API 文档按函数名和类型检索。
在仓库根目录运行
\begin{lstlisting}[language=bash,numbers=none]
bash scripts/build_api.sh
# Open tmp/doxygen/html/index.html in a browser.
\end{lstlisting}
脚本优先使用系统 Doxygen，也支持仓库 tmp 中准备好的局部工具。
若提示找不到 Doxygen，安装后重试；公式渲染使用主机 LaTeX 工具链。

以附录 C 的 \code{stepBiasWalk} 为例，先在 API 中找到声明，再对照这些字段：
\begin{center}\small
\begin{tabular}{lp{10cm}}\toprule
字段 & 阅读时要获得的信息\\\midrule
brief & 一次标量偏置更新，说明函数解决的问题\\
param & bias 是当前偏置；density 是偏置单位除以平方根秒；dt 是实际经过的秒数；unit\_noise 是标准正态样本\\
return & 下一步偏置，单位与输入 bias 相同，调用方将它保存到下一次采样\\
details / pre & 更新公式、独立样本要求，以及时间和密度的有效范围\\\bottomrule
\end{tabular}
\end{center}
随后到源文件找到式\eqref{eq:bias_step}，用 $\Delta t=0.01$ s、$\xi=1$ 手算一次调用。
最后打开偏置演示，查看随机数如何产生、返回值如何保存。
从“公式—声明—函数体—调用点”走一遍，就能看懂一个小模块怎样接入程序。

自己编写注释时，用输入的物理含义解释单位与前提。例如
“dt：实际经过的时间，单位 s”比单写“dt：时间”更有帮助。
数学可用 Doxygen 的行内或行间公式语法书写，再生成页面检查公式是否正确显示。

\section{把实验写成别人可以重做的步骤}
可复制仓库中的报告模板：
\begin{center}\code{docs/EXPERIMENT_REPORT_TEMPLATE.md}\end{center}
先提出一个能比较的问题，例如“偏置密度加倍后，10 秒终点方差是否变为四倍”。
写出模型预测，然后列出运行版本、参数、种子、时长与重复次数，提供生成原始数据的命令。
结果中同时保留每组样本数、成功或失败状态、统计指标与单位，读者便能重算你的结论。

性能实验按第 20 章分别记录准备、计算、传输和消息阶段，注明预热与计时边界；
轨迹实验按第 16--19 章说明空间检查、动力学界、时间对齐和跟踪指标。
当结果偏离预测时，先判断是输入、坐标、时间、数值误差还是模型条件发生了变化，
再有针对性地修改下一组实验。

\section{练习：重放同一段运动}
\begin{enumerate}
\item \textbf{代码 appD-1：}从混合话题列表中选出扫描与原始 IMU；缺少必要输入时说明缺失项。
\item 对 0.2 rad/s、采集间隔 0.005 s 的陀螺数据，分别求 0.5 倍和 2 倍播放时的墙钟间隔与角度增量。
\item 运行自动回放例子，读取 summary.json 和 bag 元数据，解释采集与存储时间的差异。
\item 将一条 IMU 的协方差改动一个元素，预测字段摘要是否改变；说明为何检查时包含协方差与 frame。
\item 画出回放 TF 树，标注每条边的发布者；更换传感器安装位置时应改哪两条边？
\item 在 API 中找到偏置函数，从声明追到调用点；复制报告模板，填写附录 C 的两频率比较。
\end{enumerate}
练习入口与手算提示见 \code{exercises/appD/README.md}。
